\section{Расширенная форма Бэкуса-Наура}
\label{section:EBNF}

Контекстно-свободные грамматики позволяют компактно описывать многие языки, однако для выражения повторяющихся конструкций в них приходится прибегать исключительно к рекурсии.
Например, однозначная грамматика правильных скобочных последовательностей $S \to a S b S \mid \varepsilon$ использует рекурсивный нетерминал $S$ в собственной правой части дважды, хотя по смыслу конструкция $a S b$ должна повторяться произвольное число раз.
Расширенная форма Бэкуса-Наура (Extended Backus-Naur Form, EBNF)~\sidecite{Wirth1977} разрешает использовать конструкции регулярных выражений в правых частях правил грамматики, делая запись более компактной и наглядной.

\begin{definition}[EBNF-грамматика]
    \emph{EBNF-грамматика}~--- это четвёрка вида $\langle \Sigma, N, P, S \rangle$, где
    \begin{itemize}
        \item $\Sigma$~--- терминальный алфавит;
        \item $N$~--- нетерминальный алфавит, $\Sigma \cap N = \varnothing$;
        \item $P$~--- множество правил вида $N_i \to R$, где $N_i \in N$, а $R$~--- регулярное выражение над смешанным алфавитом $\Sigma \cup N$;
        \item $S$~--- стартовый нетерминал.
    \end{itemize}
\end{definition}

Иными словами, правая часть каждого правила есть регулярное выражение, в алфавит которого входят как терминалы, так и нетерминалы грамматики.
Как и для обычных регулярных выражений, допускаются операции конкатенации, альтернативы ($\mid$) и замыкания Клини ($^*$).
Расширенный синтаксис также включает сокращения $R^+ = R \cdot R^*$ и $R? = R \mid \varepsilon$.

Следует отметить, что термин <<EBNF>> в литературе не всегда используется единообразно.
Международный стандарт ISO/IEC 14977~\sidecite{ISO14977} определяет конкретный синтаксис расширенной формы Бэкуса-Наура, в котором, например, для разделения альтернатив используется вертикальная черта, символ $=$ заменяет стрелку, а правила завершаются точкой с запятой.
Однако в данной работе мы придерживаемся более математической нотации, близкой к традиционной записи грамматик.
Кроме того, в литературе встречаются такие термины, как <<Regular Right Part Grammar>> (RRPG) и <<Extended Context-Free Grammar>> (ECFG)~\sidecite{Hemerik2009}, обозначающие близкие, но не всегда тождественные понятия.
Поэтому при чтении литературы следует обращать внимание на конкретные определения, используемые авторами.

\begin{example}
    \label{ex:ebnf_dyck}
    Вернёмся к грамматике правильных скобочных последовательностей из раздела~\ref{CFG}.
    Вместо правила $S \to a S b S \mid \varepsilon$ в EBNF можно записать:
    \[
    S \to (a \ S \ b)^*
    \]
    Это правило в точности выражает ту же идею: конструкция $a S b$ может повторяться произвольное (в том числе нулевое) число раз.
\end{example}

Возникает естественный вопрос: увеличивает ли EBNF выразительную силу контекстно-свободных грамматик?

\begin{theorem}
    \label{thm:ebnf_cfg}
    Для любой EBNF-грамматики можно построить эквивалентную КС-грамматику.
\end{theorem}

\begin{proof}
    Покажем, как преобразовать каждое правило EBNF-грамматики в множество правил обычной КС-грамматики.
    Идея состоит в том, чтобы для регулярного выражения в правой части построить праволинейную грамматику~--- точно так же, как это делалось в разделе~\ref{RegularLanguages} при преобразовании регулярного выражения в грамматику.

    Рассмотрим правило $N \to R$, где $R$~--- регулярное выражение над $\Sigma \cup N$.
    \begin{enumerate}
        \item Построим недетерминированный конечный автомат $M_R = \langle \Sigma \cup N, Q, q_0, Q_f, \delta \rangle$, распознающий язык, задаваемый выражением $R$.
        Алгоритм построения автомата по регулярному выражению обсуждался в разделе~\ref{RegularLanguages} (например, алгоритм Томпсона или метод производных Бжозовского).
        \item Преобразуем автомат $M_R$ в праволинейную грамматику $G_R$ по стандартной процедуре (раздел~\ref{RegularLanguages}):
        для каждого перехода $(q_i, x, q_j) \in \delta$, где $x \in \Sigma \cup N$, добавляем правило $Q_i \to x\,Q_j$, а для каждого финального состояния $q_f \in Q_f$~--- правило $Q_f \to \varepsilon$.
        Здесь $Q_i, Q_j, Q_f$~--- новые нетерминалы, соответствующие состояниям автомата.
        \item Заменяем исходное правило $N \to R$ на множество правил грамматики $G_R$, при этом полагая $N = Q_0$ (нетерминал $N$ отождествляется с нетерминалом, соответствующим стартовому состоянию $q_0$).
        Если нетерминал $N$ уже имеет другие правила в грамматике (из других EBNF-правил), новые правила просто добавляются к ним.
    \end{enumerate}
    Применив данную процедуру ко всем правилам EBNF-грамматики, получим обычную КС-грамматику, задающую тот же язык.
\end{proof}

Таким образом, EBNF не увеличивает выразительную силу КС-грамматик, а лишь предоставляет более компактный синтаксис.
В некоторых случаях EBNF позволяет получить грамматику, размер которой экспоненциально меньше эквивалентной КС-грамматики, что немаловажно для практических приложений.

\begin{example}
    \label{ex:ebnf_conversion}
    Продемонстрируем преобразование EBNF-грамматики из примера~\ref{ex:ebnf_dyck} в обычную КС-грамматику.
    Исходное правило:
    \[
    S \to (a \ S \ b)^*
    \]

    Построим НКА для регулярного выражения $(a \ S \ b)^*$ над алфавитом $\{a, b, S\}$.
    Автомат содержит четыре состояния и следующие переходы:
    \[
    q_0 \xrightarrow{a} q_1 \xrightarrow{S} q_2 \xrightarrow{b} q_0
    \]
    Состояние $q_0$ является одновременно стартовым и финальным (благодаря замыканию Клини).

    Преобразуем НКА в праволинейную грамматику, сопоставив каждому состоянию новый нетерминал $Q_i$:
    \begin{align*}
        Q_0 & \to a \ Q_1 \mid \varepsilon \\
        Q_1 & \to S \ Q_2                  \\
        Q_2 & \to b \ Q_0
    \end{align*}

    Отождествим $S$ с $Q_0$ и подставим $Q_1$, $Q_2$ цепными подстановками:
    \[
    S \to a \ Q_1 \mid \varepsilon \;\Rightarrow\; S \to a \ S \ Q_2 \mid \varepsilon \;\Rightarrow\; S \to a \ S \ b \ S \mid \varepsilon
    \]
    Получили в точности однозначную грамматику правильных скобочных последовательностей из раздела~\ref{CFG}.
\end{example}
